\chapter{PENDAHULUAN}

\section{Latar Belakang}
Perkembangan teknologi kecerdasan buatan, khususnya dalam bidang \foreign{Automatic Speech Recognition} (ASR), telah mengalami lonjakan signifikan dalam satu dekade terakhir. Transisi dari pemodelan tradisional berbasis \foreign{Hidden Markov Model} (HMM) menuju arsitektur \foreign{End-to-End} (E2E) berbasis \foreign{deep learning} telah berhasil mereduksi \foreign{Word Error Rate} (WER) hingga lebih dari 50\% \cite{prabhavalkar_end--end_2023}. Penggunaan arsitektur \foreign{Transformer} telah menjadi standar \foreign{de-facto} dalam pemrosesan urutan suara karena kemampuannya dalam memodelkan dependensi jangka panjang (\foreign{long-range dependencies}) yang krusial untuk pemahaman konteks bahasa \cite{latif_transformers_2025}. Mekanisme \textit{self-attention} pada model ini memungkinkan sistem untuk memberikan bobot yang berbeda pada setiap bagian sinyal audio, sehingga meningkatkan akurasi pengenalan kata secara keseluruhan. Pemanfaatan teknologi ini diharapkan dapat memecahkan tantangan linguistik yang kompleks dalam berbagai bahasa, termasuk bahasa Indonesia.

Kebutuhan akan transkripsi otomatis semakin mendesak di berbagai sektor, termasuk lingkungan akademis dan media kreatif (siniar). Saat ini, solusi transkripsi berbasis \foreign{cloud} seperti Google \foreign{Cloud Speech-to-Text} masih menjadi pilihan populer karena kemudahan integrasinya. Namun, penelitian menunjukkan bahwa performa sistem ASR sangat bergantung pada karakteristik input data audio, di mana layanan berbayar tidak selalu memberikan hasil superior pada data percakapan alami yang kompleks dibandingkan solusi \foreign{open-source} yang dioptimasi \cite{ferraro_benchmarking_2023}. Selain itu, biaya operasional yang tinggi dan kekhawatiran terkait privasi data mendorong kebutuhan akan implementasi sistem transkripsi mandiri (\foreign{self-hosted}). Transisi menuju solusi lokal ini juga memungkinkan organisasi untuk memiliki kendali penuh atas data sensitif mereka tanpa harus memprosesnya melalui server pihak ketiga. Dengan demikian, pengadopsian solusi mandiri menjadi pilihan yang paling strategis bagi institusi yang memprioritaskan keamanan dan efisiensi jangka panjang.

Model \foreign{state-of-the-art} seperti \foreign{Whisper} yang dilatih menggunakan 680.000 jam data multibahasa secara \foreign{weakly supervised} menawarkan performa yang kompetitif bahkan dalam skenario \foreign{zero-shot transfer} \cite{radford_robust_nodate}. Kemampuan \foreign{Whisper} dalam menangani berbagai aksen dan kondisi audio yang bising menjadikannya fondasi yang kuat untuk aplikasi dunia nyata \cite{graham_evaluating_2024}. Pemanfaatan sumber daya komputasi internal berupa \foreign{Graphic Processing Unit} (GPU) pada institusi pendidikan dapat menekan biaya operasional secara signifikan dibandingkan penggunaan API komersial \cite{fathony_optimisasi_nodate}. Kemampuan model ini dalam menangani berbagai format audio yang tidak terstruktur juga memberikan fleksibilitas tambahan dalam pengolahan data suara yang beragam. Hal ini membuktikan bahwa investasi pada infrastruktur perangkat keras mandiri memberikan keuntungan kompetitif dalam penguasaan teknologi pengolahan bahasa alami.

Meskipun model dasarnya sangat kuat, implementasi \foreign{Whisper} dalam skala produksi menghadapi tantangan efisiensi dan kegunaan praktis. Varian model yang beragam, mulai dari \foreign{Tiny} hingga \foreign{Large-v3-Turbo}, memerlukan evaluasi mendalam untuk menemukan keseimbangan antara kecepatan pemrosesan (\foreign{inference latency}) dan akurasi, terutama pada data bahasa Indonesia yang memiliki variasi dialek dan intonasi yang unik \cite{isyanto_voice_2022}. Di sisi lain, akurasi model ASR murni jarang mencapai tingkat kesempurnaan 100\%, sehingga diperlukan mekanisme intervensi manusia (\textit{human-in-the-loop}) melalui antarmuka penyuntingan yang efisien. Penelitian sebelumnya menunjukkan bahwa menyediakan draf hasil transkripsi otomatis sebagai titik awal dapat secara signifikan meningkatkan produktivitas transkriptor manusia, asalkan tingkat kesalahan (\textit{Word Error Rate}) berada di bawah ambang batas tertentu \cite{gaur_effects_2016}. Sistem Skriptor dibangun menggunakan pendekatan arsitektur terdistribusi untuk memisahkan beban kerja antarmuka dari beban kerja inferensi AI yang berat. Desain ini memungkinkan aplikasi tetap responsif meskipun sedang memproses ribuan segmen audio secara simultan di latar belakang. Pemisahan tanggung jawab antara komponen sistem bertujuan untuk meningkatkan reliabilitas dan kemudahan dalam pemeliharaan kode di masa depan. Selain itu, arsitektur ini mendukung skalabilitas horisontal di mana jumlah \textit{worker} dapat ditambah sesuai dengan kebutuhan beban kerja transkripsi yang ada.

\section{Rumusan Masalah}
Berdasarkan latar belakang di atas, rumusan masalah dalam penelitian ini adalah:
\begin{enumerate}
    \item Bagaimana perbandingan performa kecepatan (\foreign{inference time}) dan efisiensi penggunaan VRAM antar varian model \foreign{Whisper} pada sistem terdistribusi?
    \item Sejauh mana tingkat akurasi (\foreign{Word Error Rate} dan \foreign{Character Error Rate}) yang dihasilkan oleh berbagai varian model \foreign{Whisper} pada data siniar bahasa Indonesia?
    \item Bagaimana rancang bangun sistem transkripsi terdistribusi menggunakan \foreign{Next.js} dan \foreign{Python Worker} dapat mendukung efisiensi alur kerja koreksi transkrip secara manual?
    \item Bagaimana efektivitas integrasi komunikasi asinkron antara aplikasi utama dan \foreign{worker} dalam menjaga stabilitas pemrosesan data audio berdurasi panjang?
\end{enumerate}

\section{Batasan Masalah}
Penelitian ini dibatasi pada hal-hal sebagai berikut:
\begin{enumerate}
    \item Model yang diuji terbatas pada varian \foreign{Whisper} (\foreign{Tiny, Small, Medium, Large-v3, Large-v3-Turbo}).
    \item Sistem dibangun menggunakan framework \foreign{Next.js} 16 dengan dukungan \foreign{TypeScript} dan arsitektur \textit{distributed task queue} berbasis \textit{Redis} untuk menjamin skalabilitas \cite{sekhar_emmanni_role_2021}.
    \item Dataset yang digunakan berasal dari korpus \foreign{Mozilla Common Voice} (Siniar \textit{Homostoria} dan \textit{Hari Minggoean}) dalam bahasa Indonesia.
    \item Infrastruktur pengujian menggunakan spesifikasi NVIDIA RTX A4000 pada lingkungan \foreign{Runpod}.
    \item Jangka waktu penelitian dan pengujian dibatasi selama lima minggu.
\end{enumerate}

\section{Tujuan Penelitian}
Tujuan dari penelitian ini adalah:
\begin{enumerate}
    \item Menganalisis perbandingan performa kecepatan dan konsumsi sumber daya dari seluruh varian model \foreign{Whisper} untuk menentukan efisiensi operasional.
    \item Mengukur tingkat akurasi (WER dan CER) model \foreign{Whisper} pada dataset bahasa Indonesia untuk mengidentifikasi varian yang paling reliabel.
    \item Merancang dan membangun arsitektur sistem transkripsi terdistribusi yang mencakup antarmuka penyuntingan interaktif untuk memudahkan koreksi hasil AI.
    \item Menguji skalabilitas dan stabilitas komunikasi antar komponen sistem dalam menangani beban kerja transkripsi dalam skala besar.
\end{enumerate}

\section{Manfaat Penelitian}
Manfaat yang diharapkan dari penelitian ini adalah:
\begin{enumerate}
    \item \textbf{Manfaat Teoritis:} Memberikan data \foreign{benchmarking} terbaru mengenai performa model ASR \foreign{Transformer} pada bahasa Indonesia serta kontribusi pada literatur rancang bangun sistem AI terdistribusi. Penelitian ini juga memperkaya kajian mengenai interaksi manusia dan komputer dalam proses koreksi data hasil otomatisasi.
    \item \textbf{Manfaat Praktis:} Menyediakan solusi sistem transkripsi mandiri (\textit{self-hosted}) bernama Skriptor yang ekonomis bagi institusi akademis dan profesional. Sistem ini membantu mengurangi biaya operasional, meningkatkan keamanan privasi data, dan mempercepat alur kerja pembuatan transkrip melalui fitur penyuntingan yang terintegrasi.
\end{enumerate}
